\section{逆三角関数}\label{節:逆三角関数}

$\sin\theta=\frac12$を満たす角は,\,$\frac\pi6$,\,$\frac{5\pi}6$など無数にある.このままでは,\,$x$から角を返す関数にはならない.\sectionref{節:逆向きの置換}で扱った枝の選択を,ここでも行う.

\begin{definition}{逆三角関数とその値の範囲}
次の区間に制限した三角関数の逆関数を,それぞれ$\arcsin$,\,$\arccos$,\,$\arctan$と書く.
\[
\begin{array}{c|c|c}
\text{逆関数}&\text{入力の範囲}&\text{返す角の範囲}\\ \hline
\arcsin x&[-1,1]&\left[-\frac\pi2,\frac\pi2\right]\\[3pt]
\arccos x&[-1,1]&[0,\pi]\\[3pt]
\arctan x&\R&\left(-\frac\pi2,\frac\pi2\right)
\end{array}
\]
例えば$\theta=\arcsin x$とは,\,$\sin\theta=x$かつ$-\frac\pi2\leqq\theta\leqq\frac\pi2$であることを表す.
\end{definition}

正弦と正接は指定した区間で狭義単調増加し,余弦は狭義単調減少するので,角は一つに決まる.正接の範囲の端を含めないのは,そこで$\cos\theta=0$になり,正接が定義されないためである.\,$\sin^{-1}x$などの表記もあるが,逆数との混同を避け,本書では$\arcsin x$を使う.

定義から$\sin(\arcsin x)=x$は$-1\leqq x\leqq1$で成り立つ.一方,\,$\arcsin(\sin\theta)=\theta$は,\,$\theta$が選んだ範囲にあるときに限る.左辺が返す値は,必ずその範囲内だからである.

\begin{example}{逆三角関数の範囲}
次の値を求めよ.
\[
\arcsin\left(\sin\frac{5\pi}6\right),\qquad
\arccos\left(\cos\frac{7\pi}6\right),\qquad
\arctan\left(\tan\frac{3\pi}4\right).
\]
\end{example}

第一の正弦は$\frac12$であり,\,$\arcsin$の値の範囲にある角は$\frac\pi6$である.第二の余弦は$-\frac{\sqrt3}2$であり,\,$[0,\pi]$の中では$\frac{5\pi}6$になる.第三の正接は$-1$であり,\,$\left(-\frac\pi2,\frac\pi2\right)$の中では$-\frac\pi4$になる.従って答えは順に
\[
\frac\pi6,\qquad\frac{5\pi}6,\qquad-\frac\pi4
\]
である.逆関数の記号を打ち消す前に,戻った角が指定された範囲にあるかを確かめる.

\subsection{逆三角関数の微分}

逆関数の微分は,もとの三角関数の微分から求められる.\,$y=\arcsin x$,\,$-1<x<1$とおく.\,$\sin y=x$の両辺を$x$で微分すると
\[
\cos y\cdot y'=1
\]
となる.ここで$-\frac\pi2<y<\frac\pi2$なので$\cos y>0$である.したがって$\cos y=\sqrt{1-\sin^2y}=\sqrt{1-x^2}$と符号まで決まり,
\[
(\arcsin x)'=\frac1{\sqrt{1-x^2}}\qquad(-1<x<1)
\]
を得る.端点では分母が$0$になり,この式は使えない.

$y=\arccos x$の場合には,\,$\cos y=x$を微分して$-\sin y\cdot y'=1$となる.\,$0<y<\pi$で$\sin y=\sqrt{1-x^2}>0$なので
\[
(\arccos x)'=-\frac1{\sqrt{1-x^2}}\qquad(-1<x<1)
\]
である.負号は,余弦が選んだ区間で減少することとも一致する.\,$y=\arctan x$では$\tan y=x$を微分して
\[
(1+\tan^2y)y'=1
\]
となり,\,$\tan y=x$を代入して
\[
(\arctan x)'=\frac1{1+x^2}\qquad(x\in\R)
\]
を得る.

合成関数の微分を使えば,\,$a>0$に対して
\begin{equation}\label{式:逆三角関数の積分}
\begin{aligned}
\int\frac1{a^2+x^2}\dd x
 &=\frac1a\arctan\frac xa+C,\\
\int\frac1{\sqrt{a^2-x^2}}\dd x
 &=\arcsin\frac xa+C\qquad(|x|<a)
\end{aligned}
\end{equation}
を得る.第1式では,\,$x$を$a$で割ることと,外側の係数$\frac1a$の両方が必要である.第2式では,根号から出る$a$と合成関数の微分の$\frac1a$が打ち消し合う.

一般の2次式が分母や根号の中にあるときは,平方完成して$(x-p)^2+a^2$や$a^2-(x-p)^2$の形を作り,\,$u=\frac{x-p}{a}$と置けば,上の公式の$a=1$の場合に帰着する.

\begin{bookpoint}
\textbf{POINT}\par
2次式を$(x-p)^2+a^2$や$a^2-(x-p)^2$と平方完成し,\,$u=\frac{x-p}{a}$と置いて$\arctan u$や$\arcsin u$の微分の形にする.
\end{bookpoint}

\begin{example}{二次式の逆数の積分}
次の不定積分を求めよ.
\[
\text{(1)}\quad\int\frac1{x^2-4x+8}\dd x,
\qquad
\text{(2)}\quad\int\frac1{\sqrt{9-4x^2}}\dd x.
\]
\end{example}

\paragraphbf{方針}
\begin{enumerate}[label=(\arabic*)]
\item 分母を$(x-2)^2+4$と平方完成し,\,$u=\frac{x-2}{2}$と置く.
\item 根号の中を$9\left\{1-\left(\frac{2x}{3}\right)^2\right\}$と書き,\,$u=\frac{2x}{3}$と置く.
\end{enumerate}

\paragraphbf{解答}
\begin{enumerate}[label=(\arabic*)]
\item $\frac12\arctan\frac{x-2}{2}+C$
\item $\frac12\arcsin\frac{2x}{3}+C$\quad$\left(-\frac32<x<\frac32\right)$
\end{enumerate}

\paragraphbf{解法}
\begin{enumerate}[label=(\arabic*)]
\item 分母は$x^2-4x+8=(x-2)^2+4$である.\,$u=\frac{x-2}{2}$と置くと,\,$(x-2)^2+4=4(1+u^2)$,\,$\dd x=2\dd u$だから
\begin{align*}
\int\frac1{x^2-4x+8}\dd x
&=\int\frac{2}{4(1+u^2)}\dd u\\
&=\frac12\arctan u+C\\
&=\frac12\arctan\frac{x-2}{2}+C.
\end{align*}
\item もとの関数が実数で定義されるのは$-\frac32<x<\frac32$である.\,$u=\frac{2x}{3}$と置くと,\,$\sqrt{9-4x^2}=3\sqrt{1-u^2}$,\,$\dd x=\frac32\dd u$だから
\begin{align*}
\int\frac1{\sqrt{9-4x^2}}\dd x
&=\int\frac{3}{2\cdot3\sqrt{1-u^2}}\dd u\\
&=\frac12\arcsin u+C\\
&=\frac12\arcsin\frac{2x}{3}+C.
\end{align*}
\end{enumerate}

\paragraphbf{三角置換との関係}
(1)は,\,\chapterref{章:置換積分}の三角置換でも解ける.\,$x-2=2\tan\theta$,\,$-\frac\pi2<\theta<\frac\pi2$と置くと,\,$\dd x=2(1+\tan^2\theta)\dd\theta$,\,$(x-2)^2+4=4(1+\tan^2\theta)$だから
\[
\int\frac1{x^2-4x+8}\dd x
=\int\frac12\dd\theta
=\frac\theta2+C
\]
となる.ここで$\theta$は$\tan\theta=\frac{x-2}{2}$を満たす$\left(-\frac\pi2,\frac\pi2\right)$の角,すなわち$\theta=\arctan\frac{x-2}{2}$である.2つの解き方は,\,$u=\tan\theta$という置換で結ばれている.逆三角関数は,三角置換で導入した角$\theta$を$x$の関数として書くための記号なのだ.
%
